\unit{강연 212 — 몫 준스킴}
\sid{212}
\src{99}
\seminar{제13년차, 1960/61, 제212호\\1961년 2월}
\paperhead{대수기하학에서의 구성 기법과 존재 정리들\\
III. 몫 준스킴}%
  {알렉상드르 그로텐디크 지음}

\subsection*{서론}
\addcontentsline{toc}{subsection}{서론}

이 강연에서 다루는 문제들은 가변 스킴에 대한 어떤 공변 함자들을
나타내려 한다는 점에서, 반변 함자들을 나타내려 했던 앞의 두 강연과
다르다. 그렇지만 몫을 취하는 과정은 강연 I과 II의 문제들을 포함하여
대수기하학의 많은 구성 문제에서 본질적이다([1], [2]). 이를테면 충실
평탄이고 준콤팩트인 사상 $T\longrightarrow S$에 관하여 $T$-준스킴
$X$ 위의 \emph{하강 데이터가 유효한가} 하는 문제는, 그 하강 데이터가
$X$ 안에 정의하는 평탄 동치관계에 대한 $X$의 몫이 존재하는가 하는
문제와 동치이다. 여기서 그 몫은 아래에서 살펴볼 합당한 성질들을
만족해야 한다. [1], A, \S~2 c에서 제기한 문제들도 아마 이 강연의
제2절에서 제기하는 문제들과 함께 해결될 것이다. 마찬가지로
$S$-스킴 $X$의 \emph{피카르 스킴}(정의는 [2], C, \S~3 참조)은
몇몇 다른 스킴들(유효 약수들의 스킴 또는 사영공간으로의 몰입 사상들의
스킴)을 평탄 동치관계로 나눈 몫으로 여러 방식으로 정의할 수 있다.
그런데 이 보조 스킴들의 정의와 구성은 기술적으로 더 간단하다. 실제로
이들은 [2], C, \S~2에서 정의한
$\underline{\operatorname{Hom}}_S(X,Y)$ 꼴의 스킴들과 그 변형들이며,
적절한 사영성 가정 아래에서 이들의 구성이 다음 강연의 주제가 될
것이다. 따라서 이 강연과 다음 강연의 결과들을 결합하면 적절한 가정
아래에서 피카르 스킴들을 구성하게 된다.

준스킴에서 몫을 취하는 문제에는 여전히 해결되지 않은 문제가 여럿
있다. 가장 중요한 문제는 제8절에서 언급한다. 현재 이 문제는
\emph{임의 종수의 곡선과 극화된 아벨 다양체 등에 대한 정수 위의
모듈라이 스킴}을 구성하는 데 남은 유일한 장애물이다. 그러므로 이
문제의 해결은 대수적 군 전문가들이 힘쓸 만한 가치가 있다.

\src{100}
\subsection*{1. 동치관계와 유효 동치관계}
\addcontentsline{toc}{subsection}{1. 동치관계와 유효 동치관계}

$\mathcal C$를 범주라 하고 $X$를 $\mathcal C$의 대상이라 하자. 사상쌍
\[
  p_1,p_2:R\rightrightarrows X
\]
을 $\mathcal C$에서의 \emph{동치쌍}이라고 하며, 그 \emph{공역}은
$X$, \emph{정의역}은 $R$이라고 한다. 이는 $\mathcal C$의 모든 대상
$T$에 대하여 대응하는 두 사상
\[
  p_1(T),p_2(T):R(T)\rightrightarrows X(T)
\]
(여기서 $\mathcal C$의 모든 대상 $Y$에 대하여
$Y(T)=\operatorname{Hom}(T,Y)$로 둔다)이 사상
\[
  R(T)\longrightarrow X(T)\times X(T)
\]
을 정의하고, 이 사상이 $R(T)$에서 집합 $X(T)$ 위의 한 동치관계의
그래프로 가는 전단사를 유도한다는 뜻이다. 공역이 $X$인 동치쌍들
사이에는 자명한 동치관계가 있다. 이 후자의 동치관계에 대한 하나의
동치류를 \emph{$X$ 안의 $\mathcal C$-동치관계}라 하며, 혼동할
염려가 없으면 간단히 동치관계라 한다.

$X\times X$가 존재하면, $X$ 안의 동치관계를 주는 것은 $X\times X$의
부분대상 $R$을 주되, $\mathcal C$의 모든 대상 $T$에 대하여
\[
  (X\times X)(T)=X(T)\times X(T)
\]
에서 $R(T)$에 대응하는 부분집합이 $X(T)$ 위의 한 동치관계의
그래프가 되도록 하는 것과 동치이다. 사영
$\operatorname{pr}_1,\operatorname{pr}_2$가 유도하는 $R$에서 $X$로
가는 사상들을 $p_1,p_2$라 하면, 앞의 조건은 $(p_1,p_2)$가
동치쌍이라는 뜻이다. $X\times X$와 섬유곱
$(R,p_2)\times_X(R,p_1)$이 존재한다는 전제 아래에서는, [2], A,
\S~1에서 설명한 일반 원리에 따라 $X(T)$ 안의 $R(T)$에 대한 집합론적
동치관계의 공리들을 $\mathcal C$ 안에서 도식으로 나타낼 수도 있다.
여기서는 이것이 필요하지 않을 것이다.

같은 정의역 $R$과 같은 공역 $X$를 갖는 사상쌍 $(p_1,p_2)$가 주어질
때마다, 그 쌍의 \emph{여핵}을 $Z$를 변수로 하는 공변 함자
\[
  \operatorname{Hom}_{p_1,p_2}(X,Z),
\]
즉 $up_1=up_2$를 만족하는 $X$에서 $Z$로 가는 사상 $u$들의 집합을
나타내는 $\mathcal C$의 대상 $Y$로 정의할 수 있다. $Y$가 존재하면
유일한 동형사상을 제외하고 결정된다. 이를 $X/(p_1,p_2)$로 쓰거나,
기호를 남용하여 $X/R$로 쓴다. 후자의 기호는 특히
\src{101}$(p_1,p_2)$가 동치쌍일 때 사용한다. 이 경우 그 쌍이 정의하는
동치관계와 $R$을 기호에서 동일시하는 것이 관례이다. $Y$를 $X$의
몫으로 본다면, 실제로 $Y$는 그 쌍 $(p_1,p_2)$가 정의하는 동치관계에만
의존함을 유의하자.

이제 사상
\[
  f:X\longrightarrow Y
\]
에서 출발하자. 이 사상으로 $X$를 ``$Y$ 위의 대상''으로 볼 수 있다.
섬유곱
\[
  R(f)=X\times_Y X
\]
가 존재한다고 가정하고, 그 두 사영을 $p_1,p_2$라 하자. 그러면
$(p_1,p_2)$는 동치쌍이며, 이를 사상 $f$에 \emph{결부된} 동치쌍이라
한다. 따라서 이 동치쌍은 하나의 동치관계를 정의하며, 이를 $f$에
결부된 동치관계라 한다.

공역이 $X$이고 공통 정의역이 $R$인 사상쌍 $(p_1,p_2)$가 다음을
만족하면 이를 \emph{유효 동치쌍}이라 한다.
\begin{enumerate}
  \item[(i)] 여핵 $X/(p_1,p_2)=Y$가 존재한다.
  \item[(ii)] 섬유곱 $X\times_Y X$가 존재한다.
  \item[(iii)] 성분이 $p_1,p_2$인 사상
  $R\longrightarrow X\times_Y X$가 동형사상이다.
\end{enumerate}
그러면 $(p_1,p_2)$는 실제로 동치쌍이다. 이 쌍이 정의하는 동치관계도
\emph{유효 동치관계}라 한다.

사상 $f:X\longrightarrow Y$가 다음을 만족하면 이를 \emph{유효
에피사상}이라 한다.
\begin{enumerate}
  \item[(i)] 섬유곱 $X\times_Y X=R$이 존재한다.
  \item[(ii)] $p_1,p_2$를 $R$에서 $X$로 가는 두 사영이라 할 때,
  몫 $X/(p_1,p_2)$가 존재한다.
  \item[(iii)] $f$가 유도하는 사상
  $X/(p_1,p_2)\longrightarrow Y$가 동형사상이다.
\end{enumerate}
그러면 $f$는 실제로 에피사상이며, 더 나아가 엄밀 에피사상이다
([1], A, \S~2.3 참조). 섬유곱 $X\times_Y X$가 존재하면 그 역도
성립한다. 에피사상 $f$가 정의하는 $X$의 몫대상도 $X$의
\emph{유효 몫}이라 한다.

앞의 정의들은 다음의 ``갈루아 대응''을 함의한다.

\src{102}
\result{명제}{1.1} $X$ 안의 유효 동치관계 $R$들의 집합과 $X$의 유효
몫 $Y$들의 집합 사이에는 자연스러운 순서들을 보존하는 일대일 대응이
있다. $R$에는 유효 몫 $X/R$이 대응하고, $Y$에는 표준 사영
$X\longrightarrow Y$가 정의하는 유효 동치관계가 대응한다. 이 관계는
두 사영을 갖춘 섬유곱 $X\times_Y X$로 정의된다.

매우 좋은 범주들(집합의 범주, 집합의 층의 범주 등)에서는 모든 몫이
유효이고 모든 동치관계가 유효이다. 주어진 준스킴 $S$ 위의 준스킴
범주와 같은 범주에서는 더 이상 그렇지 않다. 이는 $S$가 체의
스펙트럼인 경우에도, 심지어 $S$ 위에서 유한인 스킴들로 한정한
경우에도 마찬가지이다. 유효성 문제, 그리고 $S$ 위에서 유한이지 않은
준스킴의 경우에는 몫의 존재 문제조차 대부분 까다롭다.

\subsection*{2. 예: $S$ 위에서 유한인 준스킴}
\addcontentsline{toc}{subsection}{2. 예: S 위에서 유한인 준스킴}

$S$를 국소 뇌터 준스킴이라 하고, $S$ 위에서 유한인 준스킴들의
범주를 $\mathcal C$라 하자. 이 범주는 $S$ 위의 연접인 가환 대수의
층들의 범주의 반대범주와 동치이다. 또 $S$가 환 $A$의 아핀
준스킴이면, 이 범주는 $A$ 위에서 유한인 $A$-대수들(즉 $A$-가군으로서
유한 생성인 대수들)의 범주의 반대범주와도 동치이다. 따라서
$\mathcal C$에는 유한 사영극한과 유한 귀납극한이 존재함이 곧바로
따른다. 전자는 (아무런 유한성 가정 없이도\ldots) 잘 알려져 있다.
예를 들어 $S$ 위의 준스킴 $X,Y$의 섬유곱은 대응하는 대수들의
텐서곱 $B\otimes_A C$에 대응한다. 두 $A$-대수 준동형
$u,v:C\rightrightarrows B$가 정의하는 두 사상
$X\rightrightarrows Y$의 핵은 $u(c)-v(c)$들이 생성하는 아이디얼로
$B$를 나눈 몫에 대응한다. 유한 귀납극한에 관해서는 한편으로는
$A$-대수들의 보통 곱에 대응하는 유한 합들을, 다른 한편으로는
사상쌍 $X\rightrightarrows Y$의 여핵들을 생각하면 충분하다. 실제로
후자는 (곧바로 알 수 있듯이) 준동형
$u,v:C\rightrightarrows B$가 일치하는 $C$의 원소들로 이루어진
부분환에 대응한다. 이 부분환은 뇌터 가정에 의해 $A$ 위에서 유한이다.
더욱이 뇌터 가정을 사용하면, $S$ 위에서 유한인 준스킴들의 범주
$\mathcal C$에서 이렇게 구성한 유한 귀납극한들, 특히 몫들은 사실
\emph{모든} 준스킴의 범주에서도 몫임을 보일 수 있다.

\src{103}
[1]에서 말했듯이 $\mathcal C$에는 유효가 아닌 에피사상들이 있다
(또는 엄밀하지 않은 에피사상들이 있다\footnote{인쇄본에는
``엄밀하지 않은''에서 ``않은''이 빠져 있다. 지금 생각하는 범주에는
섬유곱이 존재하며, 바로 앞 문단은 유효성과 엄밀 에피사상임 사이의
동치를 확립한다.}; 섬유곱이 존재하므로 이는 같은 말이다).
평탄성 가정을 두지 않을 때 \emph{동치관계가 언제나 유효인지는 알지
못한다}. 이 방향에서 얻은 것은 가군의 형식 이론의 기본 정리를
증명하는 데 꼭 필요한 매우 부분적인 몇몇 긍정적 결과뿐이다([2], B,
정리 1 참조). 제기된 문제에서는 $S$가 대수적으로 닫힌 잉여체를 갖는
국소 아르틴환 $A$의 스펙트럼인 경우로 쉽게 귀착시킬 수 있음을
언급해 둔다. 그러나 $A$가 체인 경우에도 답은 알려져 있지 않다.

$S$ 위의 준스킴 $X$가 더 이상 $S$ 위에서 유한이라고 가정하지 않되,
$X$ 안의 동치관계 $R$에 대하여 $p_1:R\longrightarrow X$가 유한
사상인 경우도 생각할 수 있다. 이때 $R$을 유한 동치관계라 한다.
간단히 하기 위해 $S$와 $X$가 아핀이라고 가정하자. 그러면 $R$도
아핀이므로 상황은 순수한 가환대수의 상황으로 귀착된다. \emph{이
경우에도 몫 $X/R=Y$가 존재하는지, 그리고 표준 사상
$X\longrightarrow Y$가 유한인지조차 알려져 있지 않다}. (가장 간단한
경우는 $S$가 체 $k$의 스펙트럼이고 $X$가 $k[t]$의 스펙트럼, 즉 아핀
직선인 경우이다.) 물론 앞의 두 문제가 긍정적으로 해결되면, 지금의
상황에서 $R$이 유효임을 결론내릴 수 있다. 동치관계의 그래프 대신
제4절의 의미에서 $X$ 안의 동치 전그래프만 주어진 경우에도, 몫 $Y$의
\emph{존재}와 $f:X\longrightarrow Y$의 유한성 문제는 정확히 같은
방식으로 제기됨을 유의하자.

다소 임의적인 유한 동치관계로 몫을 취하는 문제는 주어진 준스킴
$X_i$들을 어떤 닫힌 부분준스킴들을 따라 ``붙여서'' 준스킴을 구성할
때 생긴다. 붙이기 법칙은 $X_i$들의 합인 준스킴 $X$ 안의 유한
동치관계로 정확히 표현된다. 또한 여기에서 제기한 문제들과 그 여러
변형을 해결하는 것은 [2]에서 시작한 방향으로 비사영적 구성의 일반
기법을 확립하기 위한 예비 조건이 될 것이라고 예상해야 한다.

필자가 아는 유일한 일반적 긍정 결과는 다음과 같다.

\result{명제}{2.1} $S$를 국소 뇌터 준스킴, $s$를 $S$의 점,
$\Omega$를 $k(s)$의 대수적으로 닫힌 확대체라 하자. 대응하는
\src{104}``섬유 함자'' $F$, 즉 $S$ 위에서 유한인 모든 $S$-스킴
$X$에 $X_s$의 $\Omega$-값점들의 집합을 대응시키는 함자를 생각하자.
이 함자는 (자명하게 왼쪽 완전이며) \emph{오른쪽 완전}이다. 즉 유한
귀납극한, 특히 사상쌍의 여핵과 가환한다.

이 결과를 $S$의 모든 ``기하적 점''에 사용하면 다음이 따른다.
전사 라디시엘 사상들을 법으로 하여 추론함으로써(즉 해당 사상들의
역을 형식적으로 덧붙여서) $\mathcal C$에서 얻는 ``몫'' 범주
$\mathcal C'$은 ``기하적'' 범주, 즉 집합의 범주와 같은 ``유한적인
성격''의 성질들을 만족하는 범주이다. 특히 그 안의 모든 동치관계는
유효이다. 따라서 $R$이 $S$ 위에서 유한인 $X$ 안의 동치관계이면
표준 사상
\[
  R\longrightarrow X\times_Y X\qquad (\text{여기서 }Y=X/R)
\]
은 \emph{라디시엘이고 전사}이다. 사실 이 사상은 단사사상이므로
전사 닫힌 몰입 사상이다.

\subsection*{3. 작용군의 경우}
\addcontentsline{toc}{subsection}{3. 작용군의 경우}

다시 $\mathcal C$가 임의의 범주라고 가정하자. $G,X$를 $\mathcal C$의
대상이라 하고, $G$가 $\mathcal C$의 대상 $X$에 작용하는
$\mathcal C$-작용군이라고 가정하자. 이는([2], A, \S~1 참조) $\mathcal
C$의 모든 대상 $T$에 대하여 $G(T)$에는 군 구조가, $X(T)$에는
$G(T)$-집합 구조가 주어져 있고, $T$가 변할 때 이 구조들이 $T$에 따라
``함자적으로 변한다''는 뜻이다. $\mathcal C$에서 곱 $G\times G$와
$G\times X$가 존재하면, 그러한 구조를 다음 두 사상의 쌍으로 정의할
수도 있다.
\[
  G\times G\longrightarrow G,
  \qquad \pi:G\times X\longrightarrow X,
\]
여기에는 $\mathcal C$의 모든 대상 $T$에 대하여 집합 $G(T)$와 $X(T)$에
대응하는 연산들이 $G(T)$를 $X(T)$에 작용하는 군으로 만들어야 한다는
조건이 붙는다. 이 공리를 $\mathcal C$의 몇몇 도식의 가환성으로 옮기는
일은 쉽지만 번거로우며, 실제로 내가 아는 모든 경우에 전혀 쓸모가 없다.

$G\times X$가 존재한다고 가정하고 다음 두 사상을 생각하자.
\[
  p_1,p_2:G\times X\rightrightarrows X
\]
여기서
\src{105}
\[
  p_1=\operatorname{pr}_2,\qquad p_2=\pi.
\]
쌍 $(p_1,p_2)$가 동치쌍일 필요충분조건은 $\mathcal C$의 모든 대상
$T$에 대하여 이 쌍이 정하는 사상
\[
  G(T)\times X(T)\simeq(G\times X)(T)
  \longrightarrow X(T)\times X(T)
\]
이 단사인 것이다. 다시 말해, 군 $G(T)$가 집합 $X(T)$에
\emph{자유롭게} 작용하는 것이다. 즉 $g\in G(T)$, $x\in X(T)$에 대하여
$g\cdot x=x$이면 $g$는 군 $G(T)$의 항등원이어야 한다. 이때 $G$가
\emph{$X$에 자유롭게 작용한다}고 한다(또는 $X$가 $G$에 관한
\emph{$\mathcal C$-주공간}이라고 한다). 그러면 쌍 $(p_1,p_2)$에 결부된
동치관계를 \emph{$X$에 자유롭게 작용하는 군 $G$가 정하는 동치관계}라고
한다. $X\times X$도 존재하고 쌍 $(p_1,p_2)$가 정하는 사상
\[
  p:G\times X\longrightarrow X\times X
\]
을 생각할 때, $G$가 자유롭게 작용한다는 조건은 $p$가
\emph{단사사상}이라는 뜻이다.

물론 $G$가 $X$에 자유롭게 작용하지 않을 때에도 $G$에 의한 $X$의 몫,
즉 앞의 쌍 $(p_1,p_2)$의 여핵이 존재하는 기준을 얻고자 한다.

이 여핵은 흔히 $X/G$로 쓰며, $G$가 왼쪽에서 작용한다면 $G\backslash
X$로 쓰는 편이 낫다(앞의 표기는 $G$가 오른쪽에서 작용하는 경우를 위해
남겨 둔다). $p$에 의한 $G\times X$의 ``상''이 존재하더라도(예를 들어
이 상을 $p$가 그것을 거쳐 인수분해되는 $X\times X$의 가장 작은
부분대상으로 정의하더라도), 이를 $R$이라 하면 $R$은 대개 $X$의
동치관계가 아니다. 이때 $R$에 의한 몫(더 정확히는 두 사영
$\operatorname{pr}_i$가 유도하는 $R$에서 $X$로 가는 사상들의 쌍에 의한
몫)으로 곧바로 넘어가려 하면, 처음 쌍 $(p_1,p_2)$의 고유한 성질들을
돌이킬 수 없이 잃는다. 따라서 작용하는 $\mathcal C$-군이 정하는 쌍에
직접 적용되는 동치관계 개념의 일반화를 찾아야 한다.

\subsection*{4. 전동치관계}
\addcontentsline{toc}{subsection}{4. 전동치관계}

모든 사상이 동형사상인 범주를 \emph{준군}이라고 부름을 상기하자. 한편
범주는 두 바탕집합 $(X,R)$, 곧 \emph{대상}들의 집합과 \emph{화살표}들의
집합에 다음 구조들을 준 것으로 정의해야 한다.
\src{106}
\begin{enumerate}
\item[(i)] 다음 두 사상
\[
  p_1,p_2:R\rightrightarrows X
\]
을 각각 \emph{시점 사상}과 \emph{종점 사상}이라고 한다.

\item[(ii)] 다음 사상
\[
  \pi:(R,p_2)\times_X(R,p_1)\longrightarrow R
\]
을 \emph{합성 사상}이라고 한다.
\end{enumerate}
이 데이터는 여기서 되풀이하지 않을 잘 알려진 공리들을 만족해야 한다.
그 공리들은 몇몇 도식의 가환성과, 다른 두 도식을 가환하게 하는 사상
$D:X\to R$의 존재(그러한 사상은 반드시 유일하다)로 명시할 수 있다.
$D$는 각 대상을 그에 대응하는 항등화살표로 보내며 다음을 만족한다.
\[
  p_1D=p_2D=\operatorname{id}_X.
\]
이 범주가 준군이라는 말은 각 화살표를 역화살표로 보내는, $R$의
\emph{대칭사상}이라 불리는 사상
\[
  s:R\longrightarrow R
\]
이 존재한다는 뜻이다(그러한 사상은 반드시 유일하다). 이는 $s$,
$\Delta$와 앞의 데이터를 사용해 만든 다른 네 도식의 가환성으로 명시할
수 있으며, 그 가운데 처음 두 식은 다음과 같다.
\[
  p_1s=p_2,\qquad p_2s=p_1.
\]

이 점들을 상기하면 ([2], A, \S~1)의 일반 정의는 특히 임의의 범주
$\mathcal C$의 두 대상 $X,R$에 $\mathcal C$-\emph{범주} 구조, 각각
$\mathcal C$-\emph{준군} 구조를 준다는 말의 뜻을 보여 준다. 정의상 이는
$\mathcal C$의 모든 대상 $T$에 대하여, 대상들의 집합이 $X(T)$이고
화살표들의 집합이 $R(T)$인 집합론적 의미의 범주 구조, 각각 준군 구조를
주되, $T$가 변할 때 이 구조들이 $T$에 따라
\guillemotleft 함자적으로 변하게\guillemotright{} 하는 것이다. 따라서
\emph{시점 사상}과 \emph{종점 사상}이라 불리는 다음 두 사상이 정해지고,
\[
  p_1,p_2:R\rightrightarrows X
\]
필요한 섬유곱이 존재할 때에는 다음 사상도 정해진다.
\src{107}
\[
  \pi:(R,p_2)\times_X(R,p_1)\longrightarrow R
\]
이를 \emph{합성 사상}이라고 한다. 이 세 데이터만으로 $X,R$ 위의 범주
구조가 결정되며, 이 데이터에 부과할 공리는 다음과 같다. 모든 $T$에
대하여 $X(T),R(T)$에 대응하는 세 데이터가 이 두 집합 위에 범주 구조
(각각 준군 구조)를 정해야 한다. 필요하다면 이는 잘 정해진 사상
\[
  D:X\longrightarrow R
\]
을 포함하는 몇몇 도식의 가환성으로 나타낼 수 있고, 준군의 경우에는 잘
정해진 사상
\[
  s:R\longrightarrow R
\]
도 포함한다. 이 도식들은 ``집합론적'' 경우와 마찬가지로 명시된다. 이처럼
공리를 번거롭게 해석하는 일은 다행히 실제 계산에는 쓸모가 없다. 데이터와
공리들을 몇몇 섬유곱 사이의 사상들과 사상들의 등식으로 나타낼 수 있다는
사실의 이론적 의의는 오직 다음과 같다. 왼쪽 완전 함자
$F:\mathcal C\to\mathcal C'$가 있으면(즉 유한곱 및 섬유곱과 가환하는
함자가 있으면), 이 함자는 $\mathcal C$-범주(각각 $\mathcal C$-준군)를
$\mathcal C'$-범주(각각 $\mathcal C'$-준군)로 보낸다(단, $\mathcal
C$에서 유한곱과 섬유곱이 존재한다고 가정한다).

실제로는 사상 $p_1,p_2,\pi,D,s$를 $\mathcal C$의 알맞은 준단체 대상
또는 단체 대상에서의 \emph{단체 연산}으로 해석하는 법을 아는 것이
중요하다. 적어도 $\mathcal C$에서 섬유곱들이 존재할 때에는 그렇다.
용어를 고정하기 위해 \emph{표준 단체들의 범주 $\mathcal S$}를 다음과
같이 도입하자. 그 대상들은
\[
  \Delta_n=[0,n],\qquad\text{단 }n\geq 0
\]
꼴의 유한집합($0$부터 $n$까지의 정수 구간)이고, 사상 또는 화살표는 이
유한집합들 사이의 \emph{임의의 사상}이다. 범주 $\mathcal S$는 유한집합
사상을 사상으로 삼는 \emph{공집합 아닌} 유한집합들의 범주와 동치임을
유의하자. $\mathcal S$에서는 \emph{공집합 아닌} 유한 대상족의 쌍대곱이
명백히 존재하고, 한 대상 아래의 두 대상의 밂(섬유곱의 쌍대 연산)도
존재한다. $\mathcal S'$은 $\mathcal S$와 대상은 같지만 사상이
$\Delta_n$들 사이의 \emph{증가 사상}인 부분범주를 뜻한다. 이 범주는
전순서가 주어진 공집합 아닌 유한집합들의 범주와 동치이다. 이 범주에서는
두 대상의
\src{108}
쌍대곱이 결코 존재하지 않고, 한 대상 $C$ 아래의 두 대상 $A,B$의 밂도
일반적으로 존재하지 않는다. 예를 들어 $C=\Delta_0$, $A=B=\Delta_1$로
두고 두 구조사상 $u:C\to A$와 $v:C\to B$를 서로 같게 잡으면 된다.
그러나 다음과 같은 몇몇 경우에는 존재한다.
\[
  C=\Delta_0,\qquad A=\Delta_m,\qquad B=\Delta_n,\qquad
  u(0)=m,\qquad v(0)=0,
\]
이때
\[
  \Delta_m\amalg_{\Delta_0}\Delta_n=\Delta_{m+n}
\]
이다.

범주 $\mathcal C$의 \emph{단체 대상}(각각 \emph{준단체 대상})은 정의상
$\mathcal S$(각각 $\mathcal S'$)에서 $\mathcal C$로 가는 반변 함자
$K$이다. 따라서 단체 대상을 제한하면 준단체 대상이 되지만, 첫 개념과
둘째 개념의 본질적 차이는 $K_n=K(\Delta_n)$에 \emph{대칭 연산}이 있다는
점이다. 이 연산들은 $\Delta_n$의 $\mathcal S$에서의 자기동형군으로 보는
$n+1$개 문자 위의 대칭군 원소들을 함자 $K$로 보낸 것에 대응한다.

이제 모든 $n$에 대하여 $\Delta'_n$(각각 $\Delta''_n$)을 다음 유한범주로
두자. 그 대상들의 집합은 $\Delta_n$이고, $\Delta_n$ 위의 무차별
순서관계(각각 자연스러운 전순서)가 이 범주를 정한다(화살표들의 집합은 그
순서관계의 그래프이다). $\Delta'_n$(각각 $\Delta''_n$)은 $\mathcal
S$(각각 $\mathcal S'$)의 대상 $\Delta_n$에 함자적으로 의존함이 명백하다.
이제 $Z$가 범주이면, $\Delta_n$이 변할 때
$\operatorname{Hom}(\Delta'_n,Z)$(각각
$\operatorname{Hom}(\Delta''_n,Z)$)는 범주 $\mathcal S$(각각 $\mathcal
S'$)에서 집합들의 범주로 가는 함자, 즉 $Z$라는 범주에 \emph{결부된}
단체 집합(각각 준단체 집합)이다. 이를 각각 $Z'$와 $Z''$라 하자. 더구나
$Z'$에서 유도되는 준단체 집합에서 $Z''$로 가는 자명한 자연 준동형이
있으며, 이 준동형이 동형일 필요충분조건은 $Z$가 준군인 것이다. 이제
다음이 성립한다.

\result{PROPOSITION}{4.1} 범주들의 범주에서 준단체 집합들의 범주로 가는
함자 $Z\longmapsto Z''$는 충만충실하며, \emph{범주}들의 범주와 다음
준단체 집합들의 범주 사이의 동치를 정한다. 곧 $\mathcal S'$에서
$(\mathrm{Ens})$로 가는 반변 함자 $K$ 가운데 위에서 지정한 꼴의 밂
$A\amalg_C B$를 집합들의 섬유곱으로 보내는 것들의 범주이다. 마찬가지로
준군들의 범주에서 단체 집합들의 범주로 가는 함자 $Z\longmapsto Z'$는
충만충실하며, \emph{준군}들의 범주와
\src{109}
단체 집합들의 범주 사이의 동치를 정한다. 곧 $\mathcal S$에서
$(\mathrm{Ens})$로 가는 반변 함자 $K$ 가운데 밂을 섬유곱으로 보내는
것들의 범주이다.

따라서 어떤 구조를 다른 구조의 말로 해석할 때 요구되는 대로 ``동형을
제외하고'' 논한다면, 범주를 특수한 준단체 집합으로, 준군을 특수한 단체
집합으로 볼 수 있다. 집합론적 경우로 환원하는 통상적인 방법으로부터
다음이 따른다.

\result{COROLLAIRE}{4.2} 앞의 명제는 다음을 바꾸어도 여전히 성립한다:
여전히 성립한다. 즉 범주, 준군, 단체 집합 대신 각각
$\mathcal C$-범주, $\mathcal C$-준군, $\mathcal C$의 단체 대상을 쓰되,
$\mathcal C$에서 섬유곱이 존재한다고 가정한다.

$\mathcal C$의 범주 $(X,R,\ldots)$에 결부된 $\mathcal C$의 반단체 대상
$K$는 다음과 같이 명시할 수 있다. $K$의 성분
$K_n=K(\Delta_n)$을 $(R,p_1)$의 $X$ 위에서의 $n$번째 섬유
거듭제곱으로 해석하거나\footnote{FGA-CANON-ERR-0001. 인쇄본에는
“$(n+1)$번째 섬유 거듭제곱”이라 되어 있고, 뒤이은 공식에는
$K_0=R$이라 되어 있다. 신경의 대상 성분은 $K_0=X$이고 화살표
성분은 $K_1=R$이므로, 두 첨자 수정은 서로 연동된다.}, 더 나아가
다음 점화식으로 정의한다.
\[
  K_0=X,\qquad
  K_n=(K_{n-1},p_{n-1}^{(n-1)})\times_X(R,p_1),
\]
여기서 $p_i^{(n-1)}$ ($0\leq i\leq n-1$)는 $K_{n-1}$에서 $X$로 가는
자연 사영들이며, 이들 역시 점화적으로 정의된다. 이렇게 하면
$p_1,p_2,\pi,D,s$는 각각 $\mathcal S$의 다음 사상에 대응하는 단체
연산으로 해석된다. 곧 $\Delta_1$의 $0$번 면, $\Delta_1$의 $1$번
면, $\Delta_2$의 $(0,2)$번 면, 퇴화 $\Delta_1\to\Delta_0$,
$\Delta_1$의 대칭이다. 나머지 모든 반단체(각각 단체) 연산은 앞의
네 개(각각 다섯 개)로부터 합성과 섬유곱에 의해 형식적으로 따른다.

이제 범주의 대상 $X$에서 일반적으로 \emph{전동치관계}란, 대상들의
대상(이렇게 말해도 된다면)이 $X$인 준군을 주는 자료를 말한다. 따라서
이런 자료에는 특히 하나의 대상 $R$과 두 사상
\[
  p_1,p_2:R\rightrightarrows X
\]
이 포함된다. 그러나 동치쌍의 경우와 달리, 이 자료들만으로는 지금
고려하는 구조가 결정되지 않는다. 이 강연에서는 몫으로 넘어갈 수
있는 판정기준, 즉 쌍 $(p_1,p_2)$의 여핵을 구성할 수 있는 판정기준을
얻기 위해 이 개념을 고려한다. 따라서 문제의 진술 자체는 준군에 내재한
추가 자료를 사용하지 않는다. 그러나 뒤이을 결과들의 증명에서는
\src{110}
이 추가 자료, 특히 차원 $3$까지의 단체 연산(대칭 연산 포함)을
사용한다. 여기에는 $R$의 $X$ 위에서의 네 겹 섬유 거듭제곱이 관여한다.

$\mathcal C$의 대상 $X$에서의 동치관계는 전동치관계를 정한다. 이는
집합의 경우만 보면 충분하다. 즉 집합 $X$의 동치관계에, 대상의
집합이 $X$이고 화살표의 집합이 그 동치관계의 그래프인 준군을 대응시킨다.

$\mathcal C$-모노이드 $G$가 $\mathcal C$의 대상 $X$에 작용하면, 바탕
대상이 $R=G\times X$와 $X$인 $\mathcal C$-범주를 정한다(단,
$G\times X$가 존재한다고 가정한다). 이 범주가 $\mathcal C$-준군일
필요충분조건은 $G$가 군인 것이다. 이 역시 집합의 경우에 확인하면
충분하다. 화살표 $(g,a)$와 $(g',g\cdot a)$의 합성은 다음으로
정의한다.
\[
  (g',g\cdot a)(g,a)=(g'g,a),
\]
즉 $a,b\in X$이면 $\operatorname{Hom}(a,b)$는 정의상 $a$를 $b$로 보내는
운반자이고, 사상들은 $G$의 원소들을 곱하여 합성한다.

\medskip
\noindent\textbf{비고. ---} (4.1)과 같은 진술이 일으키는 논리적 어려움은,
고려하는 모든 대상이 하나의 고정된 “우주”(이 우주 자체도 집합이다) 안에
있다고 암묵적으로 가정하면 피할 수 있다.

\fsection{5}{유한 평탄 동치관계로 나눈 몫}

\result{정리}{5.1} $X=\operatorname{Spec}(B)$를 아핀 스킴, $R$을 $X$의
전동치관계라 하고, 그 성분 $R_1$이 아핀이라 하자. 곧
$R_1=\operatorname{Spec}(C)$라 하자. 제1사영 $p_1:R_1\to X$가 유한이고
국소 자유인 사상이라고 가정하자. 다시 말해 대응하는 환 준동형사
$p_1^*:B\to C$에 의해 $C$가 유한 생성 사영 $B$-가군이 된다고
가정한다. 두 준동형사
\[
  p_1^*,p_2^*:B\rightrightarrows C
\]
의 쌍의 핵인 $B$의 부분환을 $A$라 하자($p_1^*(b)=p_2^*(b)$를 만족하는
원소 $b$들의 집합이다). $Y=\operatorname{Spec}(A)$라 하고, $A$에서 $B$로의
포함에 의해 정의되는 사상을 $f:X\to Y$라 하자. 이 조건들 아래에서
다음이 성립한다.
\begin{enumerate}
\item[(i)] $B$는 $A$ 위에서 정수적이다. 즉 $f$는 정수적 사상이다.

\item[(ii)] 사상 $f$는 전사이고, 그 섬유들은 $X$에서 $R$로 나눈
집합론적 동치류
\src{111}
$p_2\allowbreak(p_1^{-1}(x))$이다. 또 $Y$의 위상은 $X$의 몫위상이다.

\item[(iii)] $Y$는 준스킴의 범주에서 $X$를 $R$로 나눈 몫이다.

\item[(iv)] $R$이 \emph{동치관계}에서 나왔다면 사상 $f:X\to Y$는
유한 국소 자유이고(즉 $B$는 유한 생성 사영 $A$-가군이고), 그
동치관계는 유효하다. 즉 $R_1\to X\times_Y X$는 동형사이다.
\end{enumerate}

이 정리는 유한군 $G$가 환 $B$에 자기동형사들로 작용하는 경우와 불변
부분환 $A$에 관한 잘 알려진 정리를 일반화하며, 증명도 알려진
증명과 유사하다. (iii)은 다음과 같이 더 정밀히 진술할 수 있다.

\result{따름정리}{5.2} 표준 사상 $R_1\to X\times_Y X$는 \emph{전사}이다.

$R$을 여전히 $X$의 “유한이고 국소 자유인” 전동치관계라 하되, 이제
$X$는 임의의 준스킴이라 하자. $fp_1=fp_2$를 만족하는 준스킴 $Y$와
사상 $f:X\to Y$를 찾을 수 있고, 더나아가 $Y$ 위의 환의 층의
준동형사열
\[
  \mathcal O_Y\longrightarrow f_*(\mathcal O_X)
  \rightrightarrows g_*(\mathcal O_R)
\]
이 완전하다고 가정하자(여기서 $g=fp_i$). 그러면 이 정리에서 (i)--(iv)와
같은 결론을 얻는다. 특히 (iii)에 의해 $Y$는 $X$를 $R$로 나눈
몫이고, 따라서 유일한 동형사까지 결정된다. 이 조건들 아래에서
$X$의 전동치관계 $R$을 \emph{허용 가능}하다고 한다. 이 정의로부터
다음을 얻는다.

\result{정리}{5.3} $X$를 준스킴, $R$을 $X$의 전동치관계라 하고,
$p_1:R_1\to X$가 유한 국소 자유 사상이라 하자. $R$이 허용 가능할
필요충분조건은 $X$에서 $R$로 나눈 모든 집합론적 동치류
$p_2(p_1^{-1}(x))$가 아핀 열린집합에 포함되는 것이다. 이 조건은
$X$의 모든 유한 부분집합이 아핀 열린집합에 포함되면 항상 만족하며,
예를 들어 $X$가 아핀 스킴 위의 준사영이면 그렇다.

실제로 이때 $X$에서 $R$로 나눈 모든 동치류가 \emph{$R$-불변인}
아핀 열린집합에 포함됨을 쉽게 보일 수 있고, 정리 5.1에서 얻은
조각들을 붙여 몫 $Y$를 구성한다.
\src{112}

\result{따름정리}{5.4} 이 조건이 만족하고 더나아가 $R$이 동치관계에서
나왔다고 하자. 그러면 그 동치관계는 유효하다. 즉
$R_1\to X\times_Y X$는 동형사이고, $f:X\to Y$는 유한 국소 자유
사상이다.

“강하”에 의해 곧바로 다음을 얻는다.

\result{따름정리}{5.5} (5.4)의 조건 아래서, $X$가 $Y$ 위에서 모든
곳에서 계수 $n$일 필요충분조건은 $(R_1,p_1)$이 $X$ 위에서 모든 곳에서
계수 $n$인 것이다. $X,R_1$이 $Z$-준스킴이고 $p_1,p_2$가
$Z$-사상이라 하자. 그러면 $Y$도 $Z$-준스킴이며, $X$가 $Z$
위에서 평탄할 필요충분조건은 $Y$가 $Z$ 위에서 평탄한 것이다.

요약하면 다음과 같다.

\result{스콜리움}{} 준스킴의 유한이고 국소 자유이며 전사인 사상
$f:X\to Y$를 주는 것은, 유한 국소 자유 사상 $p_1:R\to X$를
갖고 모든 류 $p_2(p_1^{-1}(x))$가 아핀 열린집합에 포함되는 동치관계
$R$를 갖춘 준스킴 $X$를 주는 것과 동치이다.

\medskip
\noindent\textbf{비고 5.6.}

\smallskip
\noindent $1^\circ$ 뇌터 가정은 필요하지 않았다.

\smallskip
\noindent $2^\circ$ 이 몫 구성 개념은 카르티에의 “비분리 강하”를 특수한
경우로 포함한다. 이는 $f_*(\mathcal O_X)$가 $\mathcal O_Y$에 대해
$p$-기저를 갖는 유한 국소 자유 사상 $f:X\to Y$를 결정하는 문제에
해당한다. 여기서 $X$는 소수 $p>0$에 의해 구조층 $\mathcal O_X$가
소멸되는 주어진 준스킴이다. 이 결과는 국소환에 어떤 정칙성 가정도
두지 않고, $X$가 체 위의 대수적 스킴이라고 가정하지 않은 채로 쉽게
표현할 수 있음에 유의하라. $X$를 특성 $p$인 체의 스펙트럼으로 취하면
야코브슨--부르바키 이론을 얻는다.

\smallskip
\noindent $3^\circ$ 가브리엘은 이전에 체 $k$ 위의 유한 가환군에 대한
몫 구성 이론에서 정리 (5.3)의 특수한 경우를 얻었다. [(7.4)와 비교하라.]

\fsection{6}{고유 평탄 동치관계로 나눈 몫}

\result{정리}{6.1} $S$를 국소 뇌터 준스킴이라 하고, $X$를 준사영
$S$-스킴이라 하며, $R$을 $X$ 안의 전동치관계라 하자. 다음을 가정한다.
\src{113}
\begin{enumerate}
\item[(a)] $p_1:R_1\to X$는 고유하고 평탄하다.
\item[(b)] $R_1\to X\times_S X$는 유한 사상이다. 또는 (a)에 따라 이와
동치이게, 모든 섬유가 유한하다(이 조건은 $R$이 동치관계에서 오는 경우
자동으로 성립한다).
\end{enumerate}
이 조건들 아래에서 다음이 성립한다.
\begin{enumerate}
\item[(i)] $X/R=Y$가 존재하고, ($S$가 뇌터이고 $R$이 동치관계에서 오는
경우) $S$ 위에서 준사영이다.\footnote{인쇄본에는 두 번째 가정이 없다.
강연 212에 딸린 정오표는 이때 준사영성이 $R$이 동치관계에서 오는 경우에만
증명된 것으로 보인다고 명시한다.}

\item[(ii)] 표준 사상 $f:X\to Y$는 전사이고 고유이며 열린 사상이다. 그
섬유들은 $R$에 대한 $X$의 동치류 $p_2(p_1^{-1}(x))$이므로, $Y$는
$R$이 정하는 집합론적 동치관계에 의한 $X$의 몫 위상공간과 동일시된다.
끝으로 $R_1\to X\times_Y X$는 전사이다.

\item[(iii)] $R$이 동치관계에서 오는 경우, 그 동치관계는 유효하다. 즉,
$R_1\to X\times_Y X$는 동형사상이다. 또한 $f:X\to Y$는 평탄하다
(따라서 충실 평탄하다).
\end{enumerate}

증명에서는 $R$에 대한 $X$의 적절한 준단면들을 사용하여 (5.1)로
귀착한다. 증명은 슈발레 세미나에서 몫 대수군을 구성하는 것과 유사하다.

요약하면 다음과 같다.

\result{스콜리움}{} $S$가 국소 뇌터이고 $X$가 $S$ 위에서 준사영이라고
하자. $X$에서 $S$-준스킴 $Y$로 가는 고유하고 충실 평탄이며 전사인 사상
$f:X\to Y$를 주는 것은, $p_1:R\to X$가 고유하고 평탄한 $X$ 위의
동치관계 $R$을 주는 것과 동치이다.

같은 방법으로 다음 결과를 얻는다.

\result{정리}{6.2} $S$를 뇌터 준스킴, $X$를 $S$ 위에서 유한형인
준스킴, $R$을 $S$-준스킴 $X$ 안의 전동치관계라 하자. 다음을 가정한다.
\begin{enumerate}
\item[(a)] $p_1:R_1\to X$는 평탄하고 유한형이다.
\item[(b)] 사상 $R_1\to X\times_S X$는 준유한이다(즉, 모든 섬유가
유한하다).
\end{enumerate}
그러면 $R$에 대해 \emph{포화된} $X$의 조밀한 열린집합 $U$가 존재하여
다음을 만족한다.
\begin{enumerate}
\item[(i)] $R_U$가 $R$이 $U$에 유도하는 전동치관계이면, $U/R_U$가
존재하고 $S$ 위에서 유한형이다.

\item[(ii)] 표준 사상 $U\to U/R_U$는 전사이고 열린 사상이며, 그 섬유들은
$R_U$에 대한 집합론적 동치류이다(따라서 $U/R_U$는 $R_U$가 정하는
집합론적 동치관계에 의한 $U$의 몫 위상공간이다). 끝으로 사상
\[
  (R_U)_1\longrightarrow U\times_{U/R_U}U
\]
은 전사이다.
\src{114}

\item[(iii)] $R$이 동치관계에서 오는 경우에는, $U\to U/R_U$가 충실
평탄이고 $R_U$가 유효하도록 $U$를 택할 수 있다.
\end{enumerate}

이는 본질적으로 ``쌍유리적'' 성격의 결과이다.

\medskip
\noindent\textbf{비고 6.3.}

\smallskip
\noindent $1^\circ$ (6.1)과 (6.2)에서 가정 (b)가 불필요한지는 알지
못한다. 이 가정 때문에 군으로 몫을 취할 때 실제로는 안정자들이 모두
유한군인 경우에 한정해야 한다.

\smallskip
\noindent $2^\circ$ 평탄성 가정 없이도 (6.1), (6.2)와 비슷한 결과가
있는지 물을 수 있다. 이 방향의 반례는 하나도 알지 못한다. 반면 평탄성
가정을 유지하고 $p_1:R\to X$가 평탄하고 준유한이지만 유한하지 않은
동치관계로 한정하더라도, $X$가 아핀일 때 $R$이 유효하지 않은 일이 생길
수 있다. 나가타 다양체를 덮는 아핀 열린집합들 위에 유도되는 동치관계들을
취하면 된다(그 다양체에서는 원소가 두 개인 군이 ``허용 가능하지 않게''
작용한다).

\fsection{7}{응용}

서론에서 말했듯이, (6.1)의 가장 중요한 응용은 피카르 스킴을 구성하는
것과 여러 다른 ``모듈라이'' 문제의 해를 구성하는 것이다. 이 문제들은
나중에 다시 다루겠다.

시무라의 다음 결과에 대한 간단한 증명을 얻는다.

\result{명제}{7.1} $V$를 분수체가 $K$인 이산 값매김환이라 하고,
$A$를 $V$ 위에 정의된 아벨 스킴이라 하자. 그러면
$A\otimes_V K$의 한 몫과 아이소제니로 연결되는 모든 $K$ 위의 아벨 스킴
$B'$은 ``$V$에 대해 좋은 환원을 갖는다.'' 다시 말해 $V$ 위의 어떤 아벨
스킴 $B$에 대하여 $B\otimes_V K$와 동형이다.\footnote{인쇄본에는
$B\otimes_V A$라고 되어 있다. $B$는 $V$-스킴이고 이 주장은 그 일반섬유에
관한 것이므로, 바탕 인자는 반드시 $K$여야 한다.} (상기하듯 이러한 $B$는
본질적으로 유일하다.)

$B'$이 $A_K$를 한 부분군 스킴 $C'$로 나눈 몫이라고 가정해도 된다.
(주의: 일반적으로 $C'$은 ``축소''가 아니다. 즉 그 국소환들은 멱영원을
가질 것이다.) $A$ 안에서 $C'$의 ``폐포''인 닫힌 부분스킴 $C$, 다시 말해
$C_K$가 $C'$을 포함하도록 하는 $A$의 가장 작은 닫힌 부분스킴을 생각하자.
그러면 $C_K=C'$이고, $V$가 이산 값매김환이라는 사실로부터 $C$가 $V$ 위의
$A$의 부분군 스킴임을 쉽게 얻는다. $A$가
$\operatorname{Spec}(V)=S$ 위에서 고유하므로 $C$도 그러하고, 한편 $A$는
더 나아가 $S$ 위에서 사영이다. 따라서 (6.1)을 적용하여 $A/C=B$를
구성할 수 있고,
\src{115}
이것이 구하는 $B$이다.

마지막으로 본질적으로 잘 알려진 논증을 쓰면 (6.2)로부터 다음 결과를
얻을 수 있다.

\result{정리}{7.2} $S$를 아르틴 환의 스펙트럼이라 하고, $F$와 $G$를
$S$ 위의 유한형 군 스킴이라 하며, $u:F\to G$를 $S$ 위의 군 스킴의
준동형이라 하자. 다음을 가정한다.
\begin{enumerate}
\item[(i)] $F$는 $S$ 위에서 평탄하다.
\item[(ii)] $u$의 핵은 유한하다.
\end{enumerate}
이 조건들 아래에서 몫 스킴 $G/F$가 존재하고, 표준 사상 $G\to G/F$는
전사인 열린 사상이며, 그 섬유들은 $G$ 위의 $F$의 오른쪽 작용이 정하는
집합론적 동치류들이다. 끝으로 $u$가 단사사상이면 사상 $G\to G/F$는
평탄하고, 사상
\[
  G\times F\longrightarrow G\times_{G/F}G
\]
은 동형이다. 다시 말해 $G$는 $G/F$ 위의 주동차 공간이고 그 구조군은
오른쪽에서 작용하는 $F$, 더 정확히는 $G/F$ 위의 군 스킴으로 보는
$F\times_S(G/F)$이다(cf. [1], B, \S~6).

\result{따름정리}{7.3} 이 조건들 아래에서 $G$가 $S$ 위에서 평탄일
필요충분조건은 $G/F$가 $S$ 위에서 평탄인 것이다. 이 조건이 성립하면
$F$로 나눈 몫을 취하는 일은 $S$의 임의의 기저확대와 가환하며, $F$가
$G$의 정규 부분군 스킴이면 $G/F$에 $G$를 $F$로 나눈
\emph{몫군}의 구조를 줄 수 있다.

$S$가 체의 스펙트럼이면 모든 $S$-준스킴이 자동으로 $S$ 위에서
평탄하므로 상황은 특히 간단하다. 다음을 얻는다.

\result{따름정리}{7.4} $F,G$를 체 $k$ 위의 유한형 군 스킴이라 하고,
$u:F\to G$를 $k$-군의 준동형이라 하자. 그러면 $u$는
$F\to F'\to G$로 인수분해된다. 여기서 $F\to F'$은 $F$의 닫힌 부분군
$\operatorname{Ker}u$로 나누어 몫을 취하는 준동형이고, $F'\to G$는
닫힌 몰입 사상인 군 준동형이다. 몫 $G/F=G/F'$이 존재한다. 통상적인
형식 체계(가령 뇌터 정리류와 같은 것)가 $k$ 위의 대수군들 사이에서
성립한다.

이 결과를 쓰면 고전적 의미의 대수군(즉 $k$ 위에서 기약이고 $k$ 위에서
매끄러운 대수군)에서 몫을 취하는 일과 카르티에가 고찰한 ``무한소''
부분군으로 몫을 취하는 일을 통일적으로 다룰 수 있다.
\src{116}
뒤도네의 형식군 연구에 이어 카르티에가 도입한 ``하이퍼대수''를 $k$ 위의
형식적 스킴의 범주 안의 군으로 보고, 해당하는 경우(그것들이 $k$ 위에서
유한 계수인 하이퍼대수에 대응하는 경우)에는 $k$ 위의 유한 대수군으로
보는 것이 유리하다.

\fsection{8}{한 추측}

이는 ``힐베르트 스킴''의 몇몇 부분스킴에 작용하는 사영군으로 몫을
취할 필요에 답한다. (스킴 이론에서 힐베르트 스킴은 차우 다양체를
대신한다.)

$S$를 준스킴이라 하고 $n$을 정수라 하자. $S$ 위의 각 준스킴 $S'$에
$\mathcal O_{S'}$의 단면환에 값을 갖는 $n$행 $n$열 가역행렬들의 군
$\operatorname{Gl}(n,\Gamma(S',\mathcal O_{S'}))$를 대응시키자. 이로써
$S'$에 대한 반변 함자를 얻으며, 이 함자가 표현 가능함은 쉽게 보인다.
따라서 이 함자는 $S$ 위의 군 스킴에 대응하고, 그 군 스킴은 실제로
$S$ 위에서 아핀이며 $\operatorname{Gl}(n)_S$로 나타낸다. 이 구성은
기저변환과 양립하므로, 실제로는 모든 것이 $\operatorname{Gl}(n)$으로
나타내는 $\mathbf Z$ 위의 한 군 스킴에서 나온다.
$\operatorname{Gl}(1)$은 곱셈군이라 하며 흔히 $\mathbf G_m$으로 나타내고,
함자
$S\longmapsto\Gamma(S,\mathcal O_S)^*$, 즉 $S$ 위의 ``가역원''들의 군에
대응한다. 명백한 준동형
$\operatorname{Gl}(1)\to\operatorname{Gl}(n)$이 있고, 몫군을 쉽게
구성할 수 있다. 이를 $\operatorname{GP}(n-1)$로 나타내며
\emph{$\mathbf Z$ 위의 $n-1$차 사영군}이라 한다. 이는 $S$에 층
$\widetilde{\operatorname{Gl}(n)}_S/
\widetilde{\operatorname{Gl}(1)}_S$의 단면군을 대응시키는 함자를
표현한다. 여기서 $\widetilde{\operatorname{Gl}(n)}_S$는 $S$ 위의
$\operatorname{Gl}(n)_S$의 단면들의 싹의 층을 뜻한다. ($S$ 위의
$\operatorname{GP}(n-1)_S$의 단면들이 일반적으로 $S$ 위의
$\operatorname{Gl}(n)_S$의 단면들에서 오지는 않는다는 점에 주의해야
한다!) 또한 적어도 $S$가 뇌터일 때 $\operatorname{GP}(n-1)$은 함자
$S\longmapsto\operatorname{Aut}_S(\mathbf P_S^{n-1})$도 표현함을 증명할
수 있다. 여기서 $\mathbf P_S^{n-1}$은 $S$ 위의 상대 차원 $n-1$인 표준
사영 스킴이다. 이것이 사영군이 모듈라이 이론에 들어오는 방식이다.

$S$를 뇌터 스킴이라 하자. 원하면 아핀이라고 가정해도 된다. $X$를
준사영 $S$-준스킴이라 하고, $S$에 대해 매우 풍부한 가역층
$\mathcal L$을 갖추었다고 하자. 군 $G=\operatorname{GP}(n)_S$가 $X$에
작용하는 동시에 $\mathcal L$에도 작용하고(그 $X$ 위의 작용과 양립하는
방식으로), $X$에 \emph{자유롭게} 작용한다고 가정하자.\footnote{인쇄본에는
``$S$에''라고 되어 있다. 여기서 생각하는 작용과 몫 조건은 $X$에 관한
것이다.}
\src{117}

\result{추측}{8.1}\footnote{강연 212와 관련된 정오표는 이 추측이
``결국 거짓''이라고 지적한다. 또한 비고 8.1의 긍정적 사실이 그 뒤
증명되었다고 밝히고 멈포드의 연구와 그 정련들을 참조한다.} 앞의 조건들
아래에서 다음이 성립한다.

\smallskip
\noindent $1^\circ$ \emph{$G$가 정하는 동치관계는 유효하고, 몫
$Y=X/G$는 $S$ 위에서 유한형이며, 표준 사상 $f:X\to Y$는 평탄이고
전사이다.} (따라서 $X$는 $Y$ 위에서
$G\times_S Y=\operatorname{GP}(n)_Y$를 구조군으로 하는 동차 주다발이
된다.)

\smallskip
\noindent $2^\circ$ $\mathcal L'$을 $f$에 의한 ``충실 평탄 강하''로
$\mathcal L$에서 얻는 $Y$ 위의 가역층이라 하자(cf. [1], B,
\S~th. 1). 그러면 $\mathcal L'$은 \emph{$S$에 대해 $Y$ 위에서
예비풍부}이다. 즉 어떤 정수 $m$과 $Y$에서 알맞은 표준 사영 스킴
$\mathbf P_S^N$으로 가는 준유한 사상이 존재하여,
$(\mathcal L')^{\otimes m}$은 $\mathcal O_{\mathbf P_S^N}(1)$의 역상과
동형이다.

$X$가 $S$ 위에서 분리되어 있더라도 $Y$는 $S$ 위에서 분리되어 있지 않을
수 있음에 유의하라. (이는 전혀 병적이지 않은 ``모듈라이 문제''에서
나타나는 상황이다.) $1^\circ$가 성립하면 $Y$가 분리일 필요충분조건은
$G$가 정하는 동치관계의 그래프가 닫혀 있는 것, 즉
$G\times X\to X\times X$의 상이 닫혀 있는 것이다. (이때 이 사상은
닫힌 몰입 사상이다.) $Y$가 분리이면 $\mathcal L'$이 $S$에 대해 $Y$
위에서 예비풍부일 필요충분조건은 그것이 풍부한 것이다. 다시 말해
알맞은 텐서 거듭제곱이 사영 몰입 사상을 정의하는 것이다. 서론에서
언급한 모듈라이 문제들에서는 이렇게 얻는 동치관계의 그래프가 실제로
닫혀 있음을 보일 수 있다.

\medskip
\noindent\textbf{비고 8.1. ---} 논의를 구체화하기 위하여, 그리고 현재
실제에서 가장 중요한 경우이므로 $G=\operatorname{GP}(n)_S$라고
가정하였다. $G$에 둘 만한 합리적인 가정은 오히려 $G$가 도호쿠 군들
가운데 어느 하나의 $S$ 위의 ``꼴''이라는 것일 듯하다. (이 군들의
정수 위의 구성은 슈발레도 수행하였다.) 앞의 추측을 뒷받침하는 방향에서
강연자가 아는 유일한 긍정적 사실은 다음과 같다. $X$를 \emph{표수 $0$인
체 $k$ 위의 아핀 스킴이라 하고, 군 $\operatorname{Gl}(n)_k$ 또는
$\operatorname{GP}(n-1)_k$가 $X$에 자유롭게 작용한다고 하자. 그러면
$G$가 정하는 동치관계는 유효하고, 몫 $X/G$도 아핀이며, 사상
$X\to X/G$는 평탄이고 전사이다.} 그 증명은 다음 사실을 사용한다.
(현재로서는 표수 $0$에서만 증명되어 있다.) $G$가 $G$의 아핀 좌표환
$A$에 작용하게 하고 이를 $k$-벡터공간으로 보면, $G$의 자명 표현은
$G$-안정인 유한차원 $k$-벡터 부분공간의 한 합성열 안에서 단 한 번만
나타난다. $G$의 선형표현론을 체계적으로 사용하면
\src{118}
적어도 바탕체 위에서는 결국 이 추측의 증명을 얻을 수 있을 듯하다.
바탕이 더 이상 체가 아니면 강연자는 아무것도 알지 못한다.

\section*{참고문헌}
\addcontentsline{toc}{section}{참고문헌}

\begin{description}
\item[{[1]}] Grothendieck (Alexander). --- \emph{Techniques de descente et
théorèmes d'existence en géométrie algébrique, I: Généralités}, Séminaire
Bourbaki, 1959/60, 제190호, 29쪽.

\item[{[2]}] Grothendieck (Alexander). --- \emph{Techniques de descente et
théorèmes d'existence en géométrie algébrique, II: Le théorème d'existence en
théorie formelle des modules}, Séminaire Bourbaki, 1959/60, 제195호, 22쪽.

\item[{[3]}] Grothendieck (A.) 및 Dieudonné (J.). --- \emph{Éléments de
géométrie algébrique, I: Le langage des schémas}. --- Paris, Presses
universitaires de France, 1960 (Institut des hautes Études scientifiques,
Publications mathématiques, 4).
\end{description}

% FGA 정본 인쇄문 보존 장치: 원자적 통합 R1.
\par\smallskip
\begingroup\footnotesize
\noindent\textit{FGA-CANON-ERR-0039. 인쇄본의 독법은 C102에 계속 기록되어 있으며, FGA-CANON-DISP-0036이 수정된 현행 본문의 최소 수정을 규율한다.}\par
\noindent\textit{FGA-CANON-ERR-0027. 인쇄본의 독법은 C088에 계속 기록되어 있으며, FGA-CANON-DISP-0029가 수정된 현행 본문의 최소 수정을 규율한다.}\par
\noindent\textit{FGA-CANON-ERR-0001. 인쇄본의 독법은 C090에 계속 기록되어 있으며, FGA-CANON-DISP-0006이 수정된 현행 본문의 최소 수정을 규율한다.}\par
\noindent\textit{FGA-CANON-ERR-0040. 인쇄본의 독법은 C103에 계속 기록되어 있으며, FGA-CANON-DISP-0037이 수정된 현행 본문의 최소 수정을 규율한다.}\par
\noindent\textit{FGA-CANON-ERR-0041. 인쇄본의 독법은 C104에 계속 기록되어 있으며, FGA-CANON-DISP-0038이 수정된 현행 본문의 최소 수정을 규율한다.}\par
\endgroup
